\section{Related Work}
\label{sec:related}

\paragraph{Continual learning and catastrophic forgetting.}
Continual learning addresses the challenge of learning a sequence of tasks without
forgetting previously acquired knowledge~\citep{McCloskey1989,FrenchCF1999,Parisi2019}.
Early approaches include regularization-based methods (EWC~\citep{EWC2017}, SI~\citep{zenke2017continual}),
replay-based methods~\citep{Rebuffi2017iCaRL,Lopez-Paz2017GEM}, and
architectural approaches~\citep{Rusu2016,HAT2018}.
Our work targets the replay-free setting, where no past-task data is retained.

\paragraph{LoRA-based continual learning.}
Low-rank adaptation (LoRA)~\citep{hu2022lora} has been widely adopted for
parameter-efficient fine-tuning of large pretrained models.
In the continual setting, SeqLoRA maintains a single shared adapter throughout task sequences,
leading to interference between tasks.
IncLoRA~\citep{Wang_2022_CVPR} and related structural methods use per-task or
growing adapters to prevent interference at the cost of increased storage.
EWC-LoRA~\citep{zheng2025ewclora} applies Fisher-weighted regularization in the $\DeltaW$ space
to protect task-relevant directions without growing the adapter.
O-LoRA~\citep{Wang2023OLoRA} projects new-task gradients onto the null space of
previous adapters.
FR-LoRA~\citep{adhikari2025frlora} applies two-stage Fisher regularization for
multilingual LoRA continual learning but does not address flatness control.
Our work is most closely related to EWC-LoRA, but adds a principled flatness component
grounded in the near-orthogonality discovery.

\paragraph{Sharpness-aware optimization.}
SAM~\citep{foretSharpnessAwareMinimizationEfficiently2021} minimizes the worst-case loss
within a neighborhood of the parameters, finding flat minima with better generalization.
GAM~\citep{zhengGradientAwareminimizationDistributed2021} reformulates sharpness
via gradient norm perturbation, avoiding the inner-maximization solve.
Flat-LoRA~\citep{li2024flatlora} applies full-space SAM perturbations alongside LoRA adapters
to find flat regions in the full parameter space; it is a single-task fine-tuning method,
not a CL method.
LoRA-SAM~\citep{NEURIPS2024_4eb2c0ad} restricts SAM to LoRA coordinates.
In the continual learning context, flat minima have been linked to reduced
forgetting~\citep{NEURIPS2020_518a38cc,shiOvercomingCatastrophicForgetting2021}.
C-Flat~\citep{NEURIPS2024_C-Flat} proposes a plug-and-play zeroth-plus-first-order flatness
optimizer for CL, showing consistent gains across CL categories.
Unlike FS-LoRA, C-Flat does not combine flatness with Fisher regularization and
does not study the gradient interaction between the two objectives.
We build on adapter-subspace sharpness (AS$^1$)~\citep{Anonymous2026PaperB},
which restricts perturbations to the effective update space $\DeltaW = BA$,
making the flatness measure adapter-intrinsic rather than full-model.

\paragraph{Gradient geometry and orthogonality in continual learning.}
GEM~\citep{Lopez-Paz2017GEM} and A-GEM~\citep{chaudhry2019gem} project gradients onto
the feasible cone defined by past-task gradients.
OGD~\citep{farajtabar2020orthogonal} projects onto the orthogonal complement of previous
task gradients.
PCGrad~\citep{yu2020gradient} detects and resolves task-to-task gradient conflicts.
These methods study \emph{task-task} gradient orthogonality and enforce it as a hard
constraint.
FOPNG~\citep{garg2026fopng} projects gradients onto the Fisher-orthogonal complement of
past task gradients, combining natural gradient geometry with orthogonal projection.
The dual flatness-aware OGD framework of~\citet{yangRevisitingFlatnessAwareOptimization2025} combines
data and weight perturbation with gradient projection for CL, studying orthogonality
between task updates.
\emph{FS-LoRA studies a fundamentally different form of orthogonality}: not between
task-task updates (which requires a memory of past gradients or a constraint mechanism),
but between the \emph{flatness gradient component} $(\ggam - \gclean)$ and the
\emph{Fisher gradient} $\gfisher$ within a single task update.
This orthogonality is measured as an emergent property, not enforced as a constraint,
and it is the geometric fact that justifies the decoupled design.

\paragraph{Fisher information in continual learning.}
EWC~\citep{EWC2017} uses the diagonal Fisher to identify and protect important parameters.
Online EWC~\citep{schwarz2018progress} and SI~\citep{zenke2017continual} provide
online approximations.
In the LoRA setting, the Fisher must be computed in the $\DeltaW$ space for
basis-invariant protection~\citep{zheng2025ewclora}.
We extend this with trace normalization that makes the EWC hyperparameter
dimensionless and consistent across datasets with different gradient scales,
and with exponential accumulation decay that prevents older tasks from dominating
as the sequence grows.
